\subsection{One Angle Calculation, Two Geometric Tests}

Take two arbitrary nonzero vectors
\[
\mathbf u=(u_1,u_2),
\qquad
\mathbf v=(v_1,v_2).
\]
Let their angles of inclination be $\alpha$ and $\beta$. Then
\[
\mathbf u=\|\mathbf u\|(\cos\alpha,\sin\alpha),
\qquad
\mathbf v=\|\mathbf v\|(\cos\beta,\sin\beta).
\]

\begin{corebox}[One comparison controls both tests]
The directed turn from $\mathbf u$ to $\mathbf v$ is $\beta-\alpha$. Its sine
detects parallel directions, while its cosine detects perpendicular directions.
The coordinate tests will be derived from these geometric facts.
\end{corebox}

\begin{center}
\begin{tikzpicture}[scale=0.9]
    \coordinate (O) at (0,0);
    \coordinate (U) at (4.2,1.6);
    \coordinate (V) at (1.6,3.6);
    \draw[axis] (-0.4,0) -- (5.2,0) node[right] {$x$};
    \draw[subset,->] (O) -- (U) node[above right] {$\mathbf u$};
    \draw[subset,->] (O) -- (V) node[above] {$\mathbf v$};
    \draw (0.8,0) arc[start angle=0,end angle=20.9,radius=0.8];
    \node at (0.97,0.18) {$\alpha$};
    \draw (0.66,0.25) arc[start angle=20.9,end angle=66.0,radius=0.71];
    \node at (0.61,0.62) {$\theta$};
    \draw (1.35,0) arc[start angle=0,end angle=66.0,radius=1.35];
    \node at (1.02,1.08) {$\beta$};
\end{tikzpicture}
\end{center}

\subsubsection{Parallel directions arise from the sine of the angle difference}

\begin{derivationbox}[Deriving the determinant test]
Two nonzero vectors determine the same unoriented direction precisely when
\[
\beta-\alpha\equiv0\pmod{\pi},
\]
equivalently when $\sin(\beta-\alpha)=0$. Using
\[
\sin(\beta-\alpha)=\sin\beta\cos\alpha-\cos\beta\sin\alpha
\]
and substituting
\[
\cos\alpha=\frac{u_1}{\|\mathbf u\|},
\quad
\sin\alpha=\frac{u_2}{\|\mathbf u\|},
\quad
\cos\beta=\frac{v_1}{\|\mathbf v\|},
\quad
\sin\beta=\frac{v_2}{\|\mathbf v\|},
\]
we obtain
\[
\sin(\beta-\alpha)
=
\frac{u_1v_2-u_2v_1}{\|\mathbf u\|\,\|\mathbf v\|}.
\]
The denominator is nonzero, so the sine vanishes exactly when the numerator
vanishes.
\end{derivationbox}

\begin{theorembox}
For nonzero vectors in the plane,
\[
\boxed{
\mathbf u\parallel\mathbf v
\iff
u_1v_2-u_2v_1=0
}.
\]
Equivalently,
\[
\boxed{
\mathbf u\parallel\mathbf v
\iff
\mathbf u=\lambda\mathbf v
\text{ for some }\lambda\neq0
}.
\]
\end{theorembox}

\begin{proofbox}[Why the determinant condition gives a scalar multiple]
If $v_1\neq0$, set $\lambda=u_1/v_1$. The equality
$u_1v_2-u_2v_1=0$ gives $u_2=\lambda v_2$, hence
$\mathbf u=\lambda\mathbf v$. If $v_1=0$, then $v_2\neq0$ and the same argument
uses $\lambda=u_2/v_2$.
\end{proofbox}

\begin{interpretationbox}[Meaning of the determinant expression]
The quantity $u_1v_2-u_2v_1$ measures whether the two coordinate directions
span genuine area. It is zero exactly when both vectors lie on one line.
\end{interpretationbox}

\subsubsection{The cosine of the same difference reveals perpendicularity}

\begin{derivationbox}[Deriving the dot-product test]
The vectors are perpendicular precisely when
\[
\beta-\alpha\equiv\frac{\pi}{2}\pmod{\pi},
\]
equivalently when $\cos(\beta-\alpha)=0$. The identity
\[
\cos(\beta-\alpha)
=
\cos\beta\cos\alpha+\sin\beta\sin\alpha
\]
gives
\[
\cos(\beta-\alpha)
=
\frac{u_1v_1+u_2v_2}{\|\mathbf u\|\,\|\mathbf v\|}.
\]
Since the denominator is nonzero, perpendicularity is equivalent to the
vanishing of the numerator.
\end{derivationbox}

\begin{definitionbox}
For vectors in the plane, the \emph{dot product} is
\[
\boxed{
\mathbf u\cdot\mathbf v=u_1v_1+u_2v_2
}.
\]
In three-dimensional space,
\[
\boxed{
(u_1,u_2,u_3)\cdot(v_1,v_2,v_3)
=u_1v_1+u_2v_2+u_3v_3
}.
\]
\end{definitionbox}

\begin{theorembox}
For nonzero vectors,
\[
\boxed{
\mathbf u\perp\mathbf v
\iff
\mathbf u\cdot\mathbf v=0
}.
\]
More generally, if $\theta$ is the smaller angle between them, then
\[
\boxed{
\mathbf u\cdot\mathbf v
=
\|\mathbf u\|\,\|\mathbf v\|\cos\theta
}.
\]
\end{theorembox}

\begin{interpretationbox}[What the sign of the dot product says]
The dot product is positive for an acute angle, zero for a right angle, and
negative for an obtuse angle. It converts angular information into one scalar.
\end{interpretationbox}

\begin{examplebox}
Take
\[
\mathbf u=(3,2),
\qquad
\mathbf v=(-2,3).
\]
Then
\[
\mathbf u\cdot\mathbf v=3(-2)+2(3)=0.
\]
Therefore $\cos\theta=0$, so the vectors are perpendicular.
\end{examplebox}

\begin{summarybox}[Two coordinate tests from one geometric comparison]
The sine of the angle difference produces the parallelism test
$u_1v_2-u_2v_1=0$. The cosine of the same difference produces the
perpendicularity test $\mathbf u\cdot\mathbf v=0$. Neither rule is imposed
arbitrarily; both are consequences of the geometry of angle.
\end{summarybox}
